\centerline{\bf \Huge Теория}
\centerline{\bf \Huge КБШ + Седракян}

\begin{flushright}
        \scriptsize{}
\end{flushright}

\centerline{\bf \large Формулировка и доказательство КБШ.}

\textit{Неравенство Коши-Буняковского-Шварца.}
Пусть $a_1, a_2, \ldots, a_n$ и $b_1, b_2, \ldots, b_n$ - это действительные числа. Тогда справедливо неравенство:

$$
\left(a_1^2+a_2^2+\cdots+a_n^2\right)\left(b_1^2+b_2^2+\cdots+b_n^2\right) \geqslant\left(a_1 b_1+a_2 b_2+\cdots+a_n b_n\right)^2 .
$$


Равенство достигается тогда и только тогда, когда $\dfrac{a_1}{b_1}=\dfrac{a_2}{b_2}=\cdots=\dfrac{a_n}{b_n}$ или когда один из наборов состоит только из нулей (нулевой набор).

\begin{proof}

    Достаточно рассмотреть квадратный трехчлен относительно $x$

    $$
    \left(a_1 x-b_1\right)^2+\left(a_2 x-b_2\right)^2+\ldots+\left(a_n x-b_n\right)^2 =
    $$
    $$
    \left(a_1^2 + a_2^2 + \ldots + a_n^2\right)x^2 - 2\left(a_1b_1 + a_2b_2 + \ldots + a_nb_n\right)x + (b_1^2 + b_2^2 + \ldots + b_n^2)
    $$

    Так как это выражение является суммой квадратов действительных чисел, то оно всегда неотрицательно. Стало быть дискриминант (мы посчитаем чётный дискриминант) такого уравнения меньше либо равен нуля. А именно:

    $$
    \dfrac{D}{4} = \left(a_1b_1 + a_2b_2 + \ldots + a_nb_n\right)^2 - \left(a_1^2 + a_2^2 + \ldots + a_n^2\right)(b_1^2 + b_2^2 + \ldots + b_n^2) \leqslant 0
    $$
    
    А это и есть, что мы хотели доказать. Из этого доказательства, кстати, очень легко видно, почему равенство достигается именно в том случае, который описан в условии.

\end{proof}

\newpage

\hspace{3mm}

\centerline{\bf \large Примеры применения КБШ на деле.}

\textit{Пример 1.} 
Зачастую встречаются в неравенствах, где слева стоит константа, умноженная на какое-то выражение, а справа просто выражение. Например, нам нужно доказать вот такое неравенство:

$$
    n(a_1^2 + a_2^2 + \ldots a_n^2) \geqslant (a_1 + a_2 + \ldots + a_n)^2
$$

Вот эту константу $n$ удобно бывает представить как $n = 1^2 + 1^2 + \ldots + 1$ --- сумма $n$ единичек. Тогда по неравенству КБШ получаем:

$$
    \left(a_1^2+a_2^2+\cdots+a_n^2\right)\left(1^2+1^2+\cdots+1^2\right) \geqslant\left(a_1 \cdot 1+a_2 \cdot 1+\cdots+a_n \cdot 1\right)^2
$$

\textit{Пример 2.}
Если в условии уже дано что-то про сумму квадратов, это хороший повод попробовать КБШ.

Пусть $a^2+b^2+c^2=2$. Докажите неравенство $a+b+c \leqslant a b c+2$.

Поймём для начала что-то про $b c$. Перепишем равенство из условия в виде $b^2+c^2=2-a^2$. К тому же верна следующая цепочка неравенств

$$
2 b c \leqslant b^2+c^2=2-a^2 \leqslant 2
$$


Значит, получаем, что $b c \leqslant 1$. Перенесём теперь $a b c$ в левую сторону и запишем КБШ

$$
b+c+a(1-b c) \leqslant \sqrt{\left((b+c)^2+a^2\right)\left(1^2+(1-b c)^2\right)}=\sqrt{(2+2 b c)\left(1+(1-b c)^2\right)}
$$


Получаем, что нам надо доказать следующее неравенство $\sqrt{(2+2 b c)\left(1+(1-b c)^2\right)} \leqslant 2$. Возведём в квадрат, сделаем замену $1-b c=x$, где $x$ неотрицательный, и сделаем преобразования

$$
\begin{gathered}
(4-2 x)\left(1+x^2\right) \leqslant 4 \\
4 x^2-2 x-2 x^3 \leqslant 0 \\
-2 x(x-1)^2 \leqslant 0
\end{gathered}
$$


Последнее неравенство верно, поэтому получаем, что и наше исходное неравенство доказано.

\textit{Пример 3. (Геометрический)}

Если говорить русским языком, то \textbf{произведение длин векторов всегда не меньше чем их скалярное произведение причём в пространстве любой размерности.}

Даны числа $x, y, z$ такие, что

$$
4^x+\sin ^4 y+\ln ^6 z=16
$$


Докажите, что

$$
2^{x+1}+3 \sin ^2 y-6 \ln ^3 z \leqslant 28
$$


Используем неравенство КБШ в векторном виде.

Рассмотрим векторы $\vec{a}=\left(2^x ; \sin ^2 y ; \ln ^3 z\right)$ и $\vec{b}=(2 ; 3 ;-6)$. Скалярное произведение

$$
\vec{a} \cdot \vec{b}=2^{x+1}+3 \sin ^2 y-6 \ln ^3 z \leqslant|\vec{a}| \cdot|\vec{b}|
$$


Имеем

$$
\begin{gathered}
|\vec{b}|=\sqrt{4+9+36}=7 \\
|\vec{a}|=\sqrt{4^x+\sin ^4 y+\ln ^6 z}=4
\end{gathered}
$$


Тогда получаем, что

$$
2^{x+1}+3 \sin ^2 y-6 \ln ^3 z \leqslant 28
$$


\newpage

\hspace{3mm}

\centerline{\bf \large Формулировка и доказательство дробного КБШ.}

\textit{Неравенство Седракяна (дробный КБШ).}
Пусть $a_1, a_2, \ldots, a_n$ и $b_1, b_2, \ldots, b_n$ - это действительные положительные числа, тогда справедливо неравенство:

$$
\dfrac{a_1^2}{b_1}+\dfrac{a_2^2}{b_2}+\cdots+\dfrac{a_n^2}{b_n} \geqslant \dfrac{\left(a_1+a_2+\cdots+a_n\right)^2}{b_1+b_2+\cdots+b_n}
$$

Равенство приходится тогда и только тогда, когда $\dfrac{a_1}{b_1}=\dfrac{a_2}{b_2}=\cdots=\dfrac{a_n}{b_n}$.

\begin{proof}

    Домножим на знаменатель правой части
    
    $$
    \left(\dfrac{a_1^2}{b_1}+\dfrac{a_2^2}{b_2}+\cdots+\dfrac{a_n^2}{b_n}\right)\left(b_1+b_2+\cdots+b_n\right) \geqslant\left(a_1+a_2+\cdots+a_n\right)^2
    $$
    
    и применим КБШ для наоборов $\left(\dfrac{a_1^2}{b_1}, \dfrac{a_2^2}{b_2}, \ldots, \dfrac{a_n^2}{b_n}\right)$ и $\left(b_1, b_2, \ldots, b_n\right)$

    $$
    \left(\dfrac{a_1^2}{b_1}+\dfrac{a_2^2}{b_2}+\cdots+\dfrac{a_n^2}{b_n}\right)\left(b_1+b_2+\cdots+b_n\right) \geqslant\left(\sqrt{\dfrac{a_1^2}{b_1} \cdot b_1}+\sqrt{\dfrac{a_2^2}{b_2} \cdot b_2}+\cdots+\sqrt{\dfrac{a_n^2}{b_n} \cdot b_n}\right)^2=
    $$
    $$
    \left(a_1+a_2+\cdots+a_n\right)^2
    $$.

    Почему равенство достигается именно при $\dfrac{a_1}{b_1}=\dfrac{a_2}{b_2}=\cdots=\dfrac{a_n}{b_n}$ следует из условия обычного КБШ.


\end{proof}


\newpage

\hspace{3mm}

\centerline{\bf \large Примеры применения неравенства Седракяна на деле.}

Вообще говоря, это неравенство можно применять практически всегда, когда есть сумма каких-то дробей. Самое главное, чтобы числитель и знаменатель каждой дроби был положительный, ведь Седракян работает только для них.

\textit{Пример 1.}
Числа $a, b, c, d$ положительны и причём $a+b+c+d=1$. Найдите наименьшее возможное значение выражения

$$
\dfrac{a^2}{1-a}+\dfrac{b^2}{1-b}+\dfrac{c^2}{1-c}+\dfrac{d^2}{1-d}
$$

По неравенству Седракяна сразу же получаем

$$
\dfrac{a^2}{1-a}+\dfrac{b^2}{1-b}+\dfrac{c^2}{1-c}+\dfrac{d^2}{1-d} \geqslant \dfrac{(a+b+c+d)^2}{1-a+1-b+1-c+1-d}=\dfrac{1}{3}
$$

Причём равенство достигается, когда $a = b = c = d = \dfrac{1}{4}$

\textit{Пример 2.}
Про положительные числа $a, b, c$ известно, что $\dfrac{1}{a}+\dfrac{1}{b}+\dfrac{1}{c}=6$. Докажите неравенство

$$
\dfrac{a^2+b}{a+b}+\dfrac{b^2+c}{b+c}+\dfrac{c^2+a}{c+a} \geqslant \dfrac{9}{4}
$$

Запишем левую часть в следующем виде, и оценим сначала первую сумму, а потом сумму, которую вычиатем.

$$
\dfrac{a^2+\dfrac{1}{4}+b}{a+b}+\dfrac{b^2+\dfrac{1}{4}+c}{b+c}+\dfrac{c^2+\dfrac{1}{4}+a}{c+a}-\left(\dfrac{1}{4(a+b)}+\dfrac{1}{4(b+c)}+\dfrac{1}{4(a+c)}\right)
$$


Числитель этих дробей можно оценить с помощью следующего неравенства $a^2+\dfrac{1}{4} \geqslant a$, откуда получим такую оценку суммы:

$$
\dfrac{a^2+\dfrac{1}{4}+b}{a+b}+\dfrac{b^2+\dfrac{1}{4}+c}{b+c}+\dfrac{c^2+\dfrac{1}{4}+a}{c+a} \geqslant \dfrac{a+b}{a+b}+\dfrac{b+c}{b+c}+\dfrac{c+a}{c+a}=3
$$


А вычитаемую сумму запишем таким образом:

$$
\dfrac{\left(\dfrac{1}{4}+\dfrac{1}{4}\right)^2}{(a+b)}+\dfrac{\left(\dfrac{1}{4}+\dfrac{1}{4}\right)^2}{(b+c)}+\dfrac{\left(\dfrac{1}{4}+\dfrac{1}{4}\right)^2}{(a+c)}
$$


Мы можем оценить каждую из дробей по неравенству КБШ для дробей, откуда получим следующее

$$
\dfrac{\left(\dfrac{1}{4}+\dfrac{1}{4}\right)^2}{(a+b)}+\dfrac{\left(\dfrac{1}{4}+\dfrac{1}{4}\right)^2}{(b+c)}+\dfrac{\left(\dfrac{1}{4}+\dfrac{1}{4}\right)^2}{(a+c)} \leqslant\left(\dfrac{\dfrac{1}{4^2}}{a}+\dfrac{\dfrac{1}{4^2}}{b}\right)+\left(\dfrac{\dfrac{1}{4^2}}{b}+\dfrac{\dfrac{1}{4^2}}{c}\right)+\left(\dfrac{\dfrac{1}{4^2}}{c}+\dfrac{\dfrac{1}{4^2}}{a}\right)=
$$

$$
=\dfrac{1}{8} \cdot\left(\dfrac{1}{a}+\dfrac{1}{b}+\dfrac{1}{c}\right)=\dfrac{3}{4}
$$


В итоге получаем, что и требовалось доказать

$$
\begin{gathered}
\dfrac{a^2+b}{a+b}+\dfrac{b^2+c}{b+c}+\dfrac{c^2+a}{c+a}= \\
=\dfrac{a^2+\dfrac{1}{4}+b}{a+b}+\dfrac{b^2+\dfrac{1}{4}+c}{b+c}+\dfrac{c^2+\dfrac{1}{4}+a}{c+a}-\left(\dfrac{1}{4(a+b)}+\dfrac{1}{4(b+c)}+\dfrac{1}{4(a+c)}\right) \geqslant 3-\dfrac{3}{4}=\dfrac{9}{4}
\end{gathered}
$$

\textit{Пример 3.}
Для положительных $a, b, c, d$ докажите неравенство

$$
\frac{12}{a+b+c+d} \leqslant \frac{1}{a+b}+\frac{1}{a+c}+\frac{1}{a+d}+\frac{1}{b+c}+\frac{1}{b+d}+\frac{1}{c+d} \leqslant \frac{3}{4}\left(\frac{1}{a}+\frac{1}{b}+\frac{1}{c}+\frac{1}{d}\right)
$$

Воспользуемся КБШ для дробей:

$$
\begin{aligned}
\frac{1}{a+b}+\frac{1}{a+c}+\frac{1}{a+d}+\frac{1}{b+c} & +\frac{1}{b+d}+\frac{1}{c+d} \geqslant \frac{(1+1+1+1+1+1)^2}{3(a+b+c+d)}= \\
& =\frac{12}{a+b+c+d}
\end{aligned}
$$


Воспользуемся КБШ для дробей для правой части следующим образом:

$$
\begin{gathered}
\frac{3}{4}\left(\frac{1}{a}+\frac{1}{b}+\frac{1}{c}+\frac{1}{d}\right)=\frac{1}{4}\left(\frac{3}{a}+\frac{3}{b}+\frac{3}{c}+\frac{3}{d}\right)= \\
=\frac{3}{4}\left(\left(\frac{1}{a}+\frac{1}{b}\right)+\left(\frac{1}{a}+\frac{1}{c}\right)+\left(\frac{1}{a}+\frac{1}{d}\right)+\left(\frac{1}{b}+\frac{1}{c}\right)+\left(\frac{1}{b}+\frac{1}{d}\right)+\left(\frac{1}{c}+\frac{1}{d}\right)\right) \geqslant \\
\geqslant \frac{1}{4}\left(\frac{4}{a+b}+\frac{4}{a+c}+\frac{4}{a+d}+\frac{4}{b+c}+\frac{4}{b+d}+\frac{4}{c+d}\right)= \\
=\frac{1}{a+b}+\frac{1}{a+c}+\frac{1}{a+d}+\frac{1}{b+c}+\frac{1}{b+d}+\frac{1}{c+d}
\end{gathered}
$$

\hspace{3mm}

\centerline{\bf \large Когда нужно применять КБШ обычный и КБШ дробный.}

На самом то деле это понятно из условия самих неравенств. КБШ обычный применяем, когда видим хорошую сумму чисел в каких-то степенях, будь то просто сумма или сумма квадратов, кубов и т.д. В формулировке фигурируют два набора чисел и их соответствующие произведение. Если вы видите два набора и у них хорошие произведение чисел с одинаковыми номерами --- смело делайте КБШ. Так же, если у вас в условии дано, что сумма фиксирована помимо тригера к методу Штурма должен срабатывать и триггер к неравенству КБШ. В случае с дробным всё проще. Всё то же самое, что и с обычным КБШ, но ещё к этому добавляется, что его \textbf{ВСЕГДА} можно использовать для суммы \textbf{положительных} дробей. Даже если у вас в числителях нет квадратов, вы сами можете их организовать просто взяв набор из корней этих числителей. 